% Results template populated after the PS/PMMA production runs
\subsection{Material-realistic polymer Bragg onions}
Figure~\ref{fig:real_material_ps_pmma} compares dispersive PS/PMMA Bragg onions across the selected visible and near-infrared design wavelengths. Panel~a identifies the intact colour-performance optimum at each design wavelength, and panel~b quantifies the extent to which the corresponding outer radial reflector can be retained after progressive hollowing. Panel~c compares the calculated backward colours in CIELAB space, while panel~d shows the trade-off between backward-spectrum retention and forward attenuation for melanin-like hollow fills at 470 and 580~nm. Exact geometries, optical constants at the design wavelength and the helium d line, particle diameters, retention metrics and colour coordinates are reported in Table~\ref{tab:real_material_ps_pmma_si}.

% Insert exact numerical conclusions only after analyze_real_material_ps_pmma.py
% has generated Table_SI_real_material_complete.csv and evidence_map.csv.

\begin{figure}
  \centering
  \includegraphics[width=\linewidth]{Figure7_real_material_PS_PMMA.pdf}
  \caption{Material-realistic PS/PMMA Bragg-onion design. (a) Intact colour-performance optimisation as a function of radial-domain count. (b) Backward-integral retention during progressive inner-air substitution. (c) Calculated backward colours in the CIELAB $a^*b^*$ plane. (d) Trade-off between backward-integral retention and forward attenuation for an idealised melanin-like hollow fill with $n_{\mathrm{fill}}=1.75$.}
  \label{fig:real_material_ps_pmma}
\end{figure}
